% Float-only pages in the appendix start at the top of the page instead of being centred vertically.
\makeatletter
\setlength{\@fptop}{0pt}\setlength{\@fpbot}{0pt plus 1fil}
\setlength{\@dblfptop}{0pt}\setlength{\@dblfpbot}{0pt plus 1fil}
\makeatother
% Let the appendix's floats share pages with text and stack two wide floats at a page top.
\renewcommand{\topfraction}{0.9}\renewcommand{\bottomfraction}{0.8}\renewcommand{\textfraction}{0.05}
\renewcommand{\dbltopfraction}{0.95}\setcounter{dbltopnumber}{3}\setcounter{topnumber}{3}\setcounter{totalnumber}{6}
\renewcommand{\floatpagefraction}{0.85}\renewcommand{\dblfloatpagefraction}{0.85}
\section{Concurrent work}
\label{app:concurrent}

The works below appeared on arXiv in September and October 2026, while this paper was being finished.
Section~\ref{sec:related} contrasts the closest of them with \ourmethod in brief; here we describe each in more
detail. We have not compared against them experimentally. The fixed-point analyses of
\citet{loopsright2} and BeNice!~\citep{benice}, which we build on throughout, appeared in the same period.

\paragraph{Looped encoder--decoders.} AsymLoop~\citep{asymloop} builds on the CED of DeepSeek-V4.1-Flash and
loops the decoder. A causal encoder runs once and projects a global KV for each decoder layer, which stays fixed
across rounds; each layer of a weight-tied decoder reads it together with a full-length local KV recomputed from
the decoder's state in every round. From 60M to 1B parameters it improves on looping the whole stack, but it is
not compared with an encoder loop. AsymLoop does not discuss caching for decoding, but its decoder has a full-length KV of its own, recomputed every round,
so its cache grows with the number of rounds unless that KV is shared across rounds. Our decoder, like that of
DeepSeek-V4.1-Flash, keeps only sliding-window KV of its own: in our CED, a loop recomputes global KV every round
only if it sits in the encoder and contains a global-KV source. RLT~\citep{rlt} also pairs a parallel causal
encoder, whose output is projected into a global KV that the decoder layers read, with a decoder that keeps
sliding-window KV, but it recurs across tokens rather than across depth: each token's decoder input merges the
final decoder state of the previous token. It is evaluated on algorithmic tasks, and its decoder processes prompt
tokens one after another. Our loops repeat layers within each token and keep prefill parallel across tokens.

\paragraph{Sharing KV across rounds.} LPT~\citep{lpt} pretrains looped models in which only the first round
writes a full-length cache; later rounds read it together with a short window of their own KV. From 150M to 1B
parameters this sharing improves quality over a looped model that caches every round. We share the last round's
KV instead. Both policies cache the looped layers' KV once instead of once per round (LPT adds a short window per
round), and both compose with the reduction across layers that CED provides. In a CED, however, the choice does
not change prefill: the decoder's global KV is projected from the encoder output of the final round, so an
encoder loop runs every round on prompt tokens whichever round's KV the later rounds read. The LSTM-style looped
transformer of \citet{lstmloop} shares one cache across its whole recurrent core: every core block reads the KV
that the last core block writes in the final round, and a consistency loss late in training narrows the gap
between parallel training and recursive decoding. In an ablation, a cache producer placed before the loop gives
higher loss than this late producer at equal training tokens but lower loss at equal training compute; the author
does not conclude that either placement is better in general.

\paragraph{Compressing the caches of pretrained looped models.} HLA~\citep{hla} adds a compact cache to pretrained
Ouro models while keeping their weights frozen: exact KV for a window of recent tokens, and one latent per older
token that every loop reads directly. The cache per token shrinks by an order of magnitude. Like
LLA~\citep{lla}, HLA finds that reusing the last loop's KV in Ouro is unsafe for long generations, consistent with
Ouro's states not reaching a fixed point~\citep{loopsright2}.

\paragraph{Shared global KV with local history.} \citet{sharedglobal_local} studies a non-looped decoder in which a
lower stack produces one global KV that every upper layer reads, through one softmax, together with its own
sliding-window KV, the same split as in a CED decoder. It proves that prefill can rebuild the complete local caches
exactly while each upper layer processes only a suffix of the prompt that shrinks with depth. Our estimate of exact
SWA replay in \S\ref{sec:cost} already shrinks the replayed suffix pass by pass in this way; the proof also lets
queries and feed-forward layers run on a suffix shorter than the one used for keys and values, which would lower
that estimate slightly.

\begin{figure*}[t]
\centering
\safegraphic[width=\textwidth]{figures/fig_kvprompt.pdf}
\caption{\textbf{Last-round KV when serving a prompt.} The prompt is prefilled in one parallel pass, and the next
128 tokens are decoded one at a time; the plot shows their BPB with last-round KV minus with per-round KV, pooled
over PG19 test, FinePDFs and GovReport, with 95\% paired bootstrap intervals. Each panel compares the looped models
of one size and token budget (dense loop: L2--5; L4--7 in the last panel). The fixed-round CED encoder loop
(150M) is off the scale.}
\label{fig:kvprompt}
\end{figure*}

% Additional serving results (GB200 + H200 exports) as Appendix B: sections/appendix_serving.tex,
% figures from scripts/fig_serving_app.py.
\section{Additional Serving Results}
\label{app:serving-more}

This appendix collects serving measurements that Section~\ref{sec:cost} does not show. They come from the same two
setups: one $4\times\text{GB200}$ tray (TP4/EP4) and one $8\times\text{H200}$ node (TP8/EP8), both running SGLang
\texttt{da64c5cbb} with our plugin. Runs marked \emph{dummy weights} use random weights with standard deviation 0.02,
which route tokens to experts almost uniformly. In an A/B test, dummy weights decoded 15--36\% slower than the real
ones, while the ratios between models moved by 2--4\%. All other runs use the original (real) \vfour-Flash weights, unchanged. In Figure~\ref{fig:serving-main}, panel (a) uses dummy weights and
panels (b) and (c) use the real weights. Every
comparison is at equal rounds $T$, and the names are those of Section~\ref{sec:cost}. Across layers, a model either
shares global KV as \ced does or keeps \emph{per-layer KV} (every layer \Lr{2}{39} caches its own global KV and runs
over every prompt token in prefill). Across rounds, it keeps \emph{last-round KV} or \emph{per-round KV}. The
\emph{baseline} keeps per-layer and per-round KV; \emph{\ourmethod} keeps neither. The \emph{decoder loop} loops
\Lr{21}{38} of \ced instead of the encoder. Per-round KV is always emulated by shrinking the KV pool, so a model and its
per-round variant run identical compute.

\begin{figure*}[t]
\centering
\safegraphic[width=\textwidth]{figures/fig_serving_more.pdf}
\caption{\textbf{Further serving measurements.} \emph{(a)}~As Figure~\ref{fig:serving-main}(a), for ShareGPT on
one GB200 tray (dummy weights): output tokens/s per GPU against the p90 per-user decode speed
($1000/\text{p90 TPOT}$), one point per concurrency, preemption-free points only. \emph{(b, c)}~Mooncake tool-agent trace, closed loop, $T=8$, real
weights (the runs of Figure~\ref{fig:p1-pareto}c): \emph{(b)}~prefix-cache hit rate over the second half of each run;
\emph{(c)}~share of requests with TTFT $\le 10$\,s.}
\label{fig:serving-more}
\end{figure*}

\paragraph{The gain holds at interactive speeds.}
The $3.1\times$--$3.6\times$ of Section~\ref{sec:cost} compares peak throughput, which \ourmethod reaches at $c=512$.
Figures~\ref{fig:serving-main}(a, b) and~\ref{fig:serving-more}(a) plot throughput against the p90 per-user speed
instead. Reading the curves by linear interpolation between neighbouring points, at 30 and 50 tokens/s per user \ourmethod serves $1.39\times$--$1.67\times$ the throughput of
the baseline on 8K/1K and $1.42\times$--$1.61\times$ on ShareGPT. No KV pool is full at these speeds, so this part of
the gain comes from cross-layer sharing, not from capacity: the baseline runs every layer on every prompt token
(166 instead of 147 layer passes at $T=8$) and keeps a KV cache and an indexer in all 38 layers. A single request
shows the same: its time to first token (TTFT) on 8K/1K is 0.98\,s instead of 1.47\,s at $T=4$ and 1.68\,s instead of
2.65\,s at $T=8$, and its time per output token is 10\% lower. On H200 with real weights the gap is larger. At equal
concurrency, \ourmethod serves $1.48\times$--$1.64\times$ the throughput of the baseline on 8K/1K ($c=1$--64) and
$1.47\times$--$1.73\times$ on ShareGPT ($c=1$--256); its single-request TTFT is 1.95\,s instead of 3.23\,s ($T=4$)
and 3.56\,s instead of 5.61\,s ($T=8$). On 8K/1K the H200 baseline preempts at $c=112$ ($T=4$) and fills 96\% of its
pool at $c=64$ ($T=8$), while \ourmethod runs $c=256$ at 16\% KV use. On H200 ShareGPT the baseline's pool never
exceeds 28\%, so the $1.72\times$--$1.73\times$ there at $c=256$ comes from cross-layer sharing alone.

\begin{table*}[t]
\caption{\textbf{Cross-layer and cross-round sharing, measured separately} ($4\times\text{GB200}$, dummy weights).
Left: peak preemption-free output tokens/s per GPU on 8K/1K, with the largest preemption-free concurrency in
parentheses. Right: decode tokens/s per GPU at the largest batch of 128K-token requests that fits the KV pool, with
that batch in parentheses. Without a loop, \ced reaches 887.8 and 780.3 and the per-layer-KV model 371.7 and 116.7.
The ratios compare rows at equal $T$.}
\label{tab:serving-2x2}
\vskip 0.06in
\centering
\footnotesize
\setlength{\tabcolsep}{5pt}
\begin{tabular}{@{}llcccc@{}}
\toprule
& & \multicolumn{2}{c}{8K in / 1K out, peak} & \multicolumn{2}{c}{128K decode at the KV limit} \\
\cmidrule(lr){3-4}\cmidrule(lr){5-6}
KV across layers & KV across rounds & $T=4$ & $T=8$ & $T=4$ & $T=8$ \\
\midrule
\ced (\ourmethod, ours) & last round & 391.8 ($c$\,512) & 185.1 ($c$\,512) & 347.4 (38) & 199.2 (38) \\
\ced & per round & 308.4 ($c$\,128) & 142.5 ($c$\,64) & 186.1 (13) & 69.4 (7) \\
per layer & last round & 161.6 ($c$\,32) & 92.4 ($c$\,32) & 51.7 (3) & 29.6 (3) \\
per layer (baseline) & per round & 127.3 ($c$\,16) & 50.9 ($c$\,8) & 21.1 (1) & does not fit (0) \\
\midrule
\multicolumn{2}{@{}l}{Cross-layer sharing, both last-round KV (\ced $\div$ per layer)} & $2.43\times$ & $2.00\times$ & $6.72\times$ & $6.73\times$ \\
\multicolumn{2}{@{}l}{Cross-layer sharing, both per-round KV} & $2.42\times$ & $2.80\times$ & $8.82\times$ & --- \\
\multicolumn{2}{@{}l}{Cross-round sharing in \ced (last round $\div$ per round)} & $1.27\times$ & $1.30\times$ & $1.87\times$ & $2.87\times$ \\
\multicolumn{2}{@{}l}{Cross-round sharing with per-layer KV} & $1.27\times$ & $1.82\times$ & $2.45\times$ & --- \\
\multicolumn{2}{@{}l}{Both (\ourmethod $\div$ baseline)} & $\mathbf{3.08\times}$ & $\mathbf{3.64\times}$ & $\mathbf{16.5\times}$ & 199.2 vs.\ none \\
\bottomrule
\end{tabular}
\end{table*}

\paragraph{Both kinds of sharing contribute, and they multiply.}
Table~\ref{tab:serving-2x2} measures all four combinations of cross-layer sharing (\ced or per-layer KV) and
cross-round sharing (last-round or per-round KV). With last-round KV in both, cross-layer sharing alone gives
$2.0\times$--$2.4\times$ the peak 8K/1K throughput and $6.7\times$ the 128K decode throughput, because \ced fits 38
requests of 128K tokens where per-layer KV fits 3. Cross-round sharing alone gives $1.27\times$--$1.30\times$ in \ced
and $1.27\times$--$1.82\times$ with per-layer KV on 8K/1K, and $1.87\times$--$2.87\times$ in \ced at 128K. The 8K/1K
factors multiply to the headline exactly ($1.27\times 2.43=3.08$, $1.82\times 2.00=3.64$). Cross-round sharing thus
also helps a model with per-layer KV: the benefit of stochastic-depth training is not specific to \ced. At equal batch,
a model and its per-round variant decode within 3\% of each other (one point differs by 6\%), as the emulation intends,
so these gains come from the larger batch that fits. For the same reason, cross-round sharing changes nothing in \ced
at 30 or 50 tokens/s per user (factor 1.00), where its pool is at most a third full. It pays off once the pool fills:
at long context, on agent traces and under heavy load.

\begin{table*}[t]
\caption{\textbf{Prefill cost of running every layer over the prompt.} Top: ratio of per-request prefill time for
uncached prompts (batch 1, real weights); both models keep last-round KV. \emph{Full replay} runs the \ced decoder over
every prompt token instead of the last 128. Bottom: ratio of new prompt tokens prefilled per second behind one cached
prefix ($4\times\text{GB200}$, dummy weights, 4K new tokens per request). Values above 1 mean the \ced variant is
faster.}
\label{tab:serving-prefill}
\vskip 0.06in
\centering
\footnotesize
\setlength{\tabcolsep}{5pt}
\begin{tabular}{@{}llccccc@{}}
\toprule
& & \multicolumn{5}{c}{Prompt length (tokens)} \\
\cmidrule(lr){3-7}
& Ratio & 512 & 2K & 4K & 16K & 64K \\
\midrule
H200 & \ced, no loop: full replay $\div$ 128-token tail & 1.32 & 1.45 & 1.50 & 1.60 & 1.75 \\
H200 & Per-layer KV $\div$ \ced, encoder loop $T=8$ & 1.87 & 1.45 & 1.44 & 1.69 & 2.43 \\
GB200 & Per-layer KV $\div$ \ced, encoder loop $T=8$ & 1.65 & 1.48 & 1.41 & 1.40 & 1.48 \\
\midrule
& & \multicolumn{5}{c}{Context length / prefill chunk (tokens)} \\
\cmidrule(lr){3-7}
& & 64K / 4K & 64K / 8K & 64K / 16K & 128K / 4K & 128K / 8K \\
\midrule
GB200 & \ourmethod $\div$ baseline, $T=4$ & 1.60 & 1.64 & 1.70 & 1.65 & 1.83 \\
\bottomrule
\end{tabular}
\end{table*}

\paragraph{Prefill: skipping the decoder survives looping.}
A \ced prompt runs the encoder and a 128-token decoder tail, and an encoder loop keeps that property
(Table~\ref{tab:serving-prefill}). Even without a loop, running the decoder over the whole prompt costs $1.32\times$ at
512 tokens and $1.75\times$ at 64K tokens on H200, and $1.30\times$ at 8K tokens on GB200. With eight rounds, the
model with per-layer KV takes $1.40\times$--$1.65\times$ (GB200) and $1.44\times$--$2.43\times$ (H200) as long as the
\ced encoder loop to prefill the same prompt. Behind a cached 60K--124K-token prefix, \ourmethod at $T=4$ prefills
$1.60\times$--$1.83\times$ as many new tokens per second as the baseline at the same chunk size, and $1.56\times$ at
$T=8$ (64K context, 4K chunk). The encoder loop's own prefill grows more slowly than its round count
(Table~\ref{tab:unit-costs-app}).

\begin{table*}[t]
\caption{\textbf{Requests that fit the KV pool at long context.} $4\times\text{GB200}$: one tray with the flags of
Figure~\ref{fig:p1-pareto}, dummy weights; requests of $L+256$ tokens; the last column is decode tokens/s per GPU at
the largest 128K batch that fits. The \ced limit of 604 at 8K is computed from the pool (the sweep stopped at 256).
$8\times\text{H200}$: one node, real weights, requests that fit at full length. The H200 per-round pool also charges
each round's sliding-window rings; the GB200 one does not.}
\label{tab:serving-longctx}
\vskip 0.06in
\centering
\footnotesize
\setlength{\tabcolsep}{5pt}
\begin{tabular}{@{}llcccccc@{}}
\toprule
& & \multicolumn{3}{c}{$4\times\text{GB200}$} & \multicolumn{3}{c}{$8\times\text{H200}$} \\
\cmidrule(lr){3-5}\cmidrule(lr){6-8}
Model & KV across rounds & 8K & 128K & tok/s/GPU at 128K & 120K & 524K & 1M \\
\midrule
\ced, no loop & --- & 604 & 38 & 780.3 & 195 & 44 & 23 \\
\ced encoder loop, $T=4$ & last round & 604 & 38 & 347.4 & --- & --- & --- \\
\ced encoder loop, $T=8$ & last round & 604 & 38 & 199.2 & 195 & 44 & 23 \\
\ced decoder loop, $T=8$ & last round & --- & --- & --- & 195 & 44 & 23 \\
\ced encoder loop, $T=4$ & per round & 215 & 13 & 186.1 & --- & --- & --- \\
\ced encoder loop, $T=8$ & per round & 116 & 7 & 69.4 & 23 & 5 & 2 \\
\midrule
Per-layer KV, no loop & --- & 49 & 3 & 116.7 & 16 & 3 & 1 \\
Per-layer KV, $T=4$ & last round & 49 & 3 & 51.7 & --- & --- & --- \\
Per-layer KV, $T=8$ & last round & 49 & 3 & 29.6 & 16 & 3 & 1 \\
Per-layer KV, $T=4$ & per round & 25 & 1 & 21.1 & --- & --- & --- \\
Per-layer KV, $T=8$ & per round & 15 & 0 & --- & --- & --- & --- \\
\bottomrule
\end{tabular}
\end{table*}

\paragraph{Cost per request on four workloads.}
Figure~\ref{fig:serving-main}(c) runs four closed-loop workloads on one GB200 tray with real weights at $T=8$, one run
each:
synthetic reasoning (2K-token prompts, 4K output tokens, 32 concurrent requests), the Mooncake conversation and
tool-agent traces~\citep{mooncake} (64 concurrent requests, starting at minute 10), and synthetic long documents
(100K-token prompts, 1K output tokens, 8 concurrent requests). \ourmethod costs $3.7\times$--$4.7\times$ the GPU time
per request of an unlooped \ced. That is $0.64\times$--$0.93\times$ the cost of per-layer KV with last-round KV, and
$0.36\times$--$0.93\times$ the cost of the baseline. The largest saving is on long documents, where the baseline's pool
holds one 101K-token request at a time. The cost model, evaluated for the same runs, matches the two \ced loops
(measured over predicted 0.78--1.04). It underpredicts per-layer KV with last-round KV by 9--44\% on the three workloads
where that model's KV pool fills (95--100\%), an effect that a model of FLOPs and bytes does not capture. The encoder
loop costs the same as the decoder loop on reasoning ($1.00\times$) and $1.37\times$--$1.64\times$ as much on the three
prompt-heavy workloads, where its rounds run over every prompt token. These runs decode 8--64 requests at a time on one tray. With larger decode batches, as in
expert-parallel serving across many GPUs, each output token costs less and prompt tokens take a larger share of a
request, so the premium grows; running prompts at fewer rounds ($T_p<T_d$, Appendix~\ref{app:split}) lowers it.

\paragraph{Cross-round sharing leaves prefix caching as it is without a loop.}
Figure~\ref{fig:serving-more}(b, c) replays the Mooncake tool-agent trace in closed loop
with 64--512 concurrent sessions, real weights and $T=8$ (the runs of Figure~\ref{fig:p1-pareto}c); each point replays
600--1{,}536 complete requests from the trace. With last-round KV, both the encoder and
the decoder loop hit the prefix cache within 0.8 points as often as the unlooped \ced at every concurrency, so looping
does not change which session histories stay resident. With per-round KV, the 1.5M-token pool is full from $c=128$,
and the hit rate is 4.3, 6.4, 8.1 and 11.6 points lower than with last-round KV at $c=64$, 128, 256 and 512. The
share of requests with TTFT $\le 10$\,s then falls to 4.8\% at $c=256$ (58.3\% with last-round KV) and 1.4\% at
$c=512$ (46.0\%). A synthetic multi-turn agent workload on H200 (80 sessions of 4 turns, 32K--40K-token contexts,
0.2 requests/s) shows the same direction: per-round KV lowers the hit rate of the encoder loop from 72.3\% to 63.8\%
and its completed rate from 0.199 to 0.171 requests/s. Cross-layer sharing matters even more on this workload. In the
same runs, per-layer KV with last-round KV holds 2.0M tokens against \ced's 23.4M, so the session histories (about
3.2M tokens) do not stay resident: its hit rate is 3.4\% and it completes 0.039 requests/s, $5.1\times$ fewer than
\ourmethod.

\paragraph{Encoder and decoder loops decode at the same cost.}
Decode runs every layer under either loop placement. In a pure-decode sweep with real weights on GB200 that fills the KV
pool (8K, 16K and 32K-token contexts, 256 output tokens per request), the encoder loop and the decoder loop at $T=8$
reach 1{,}451 / 1{,}118 / 818 and 1{,}460 / 1{,}112 / 828 output tokens/s per GPU (the 8K values are capped at 768
requests). With last-round KV neither loop adds global KV, so both hold the same pool, and both take 33.1\,ms per decode
step at 8K context and batch 32.

\paragraph{Long and very long contexts.}
With last-round KV, a looped \ced fits as many requests as the unlooped \ced at every $T$
(Table~\ref{tab:serving-longctx}): 38 requests of 128K tokens on one GB200 tray, and 44 of 524K tokens and 23 of 1M
tokens on one H200 node, where the unlooped model with per-layer KV fits 3 and 1. At 524K tokens on H200, last-round KV
decodes 195 tokens/s per node at a batch of 32. Per-round KV fits only 5 such requests and peaks at 55.7 tokens/s at a
batch of 4, $3.5\times$ less. For every model, the long-context decode measurements use KV that is allocated without
being computed, so that they measure decode alone; on H200, decoding over such KV was within 3\% of decoding after a
real prefill.


\section{Evaluation protocol}
\label{app:eval}

\paragraph{Validation BPB.} We tokenize the FineWeb-Edu validation split into one token stream and cut it into
consecutive samples of \val{cfg/seq} tokens, the training context; a sample may span document boundaries. Every
model is scored on all \val{cfg/fwwin} samples, about \val{cfg/fwtok}M tokens. BPB is the total negative
log-likelihood of the predicted tokens in bits, divided by their total number of UTF-8 bytes, so it does not depend
on the tokenizer. Looped models run \val{cfg/rounds} rounds. During training we evaluated only the first
\val{cfg/evalwin/x40} samples (\val{cfg/evalwin/long} for the \val{cfg/tokens/long}B-token runs; BPB$_{\text{train}}$ in
Table~\ref{tab:long}); on the first \val{cfg/evalwin/x40} samples, every model's BPB is within 0.001 of its value on
all samples, and no ranking of models changes.

\paragraph{Last-round KV.} Last-round KV cannot be scored in one parallel pass, because a token's last-round KV
exists only after that token has run all its rounds. We therefore decode every sample token by token with teacher
forcing, in fp32, as \citet{benice} do for shared-KV perplexity. With per-round KV, the parallel pass gives the
same result as token-by-token decoding, which we check in our tests. The deployment protocol of \S\ref{sec:res-kv}
instead prefills a prompt in one parallel pass and decodes only the answer. With last-round KV the prompt keeps only its tokens' final-round KV,
which every round of the answer reads; Figure~\ref{fig:kvprompt} shows that this loses nothing at any prompt length. Table~\ref{tab:kvshare} decodes the
first \val{cfg/evalwin/x40} validation samples. For the CED encoder loop we also decode all
\val{cfg/kvfull/windows} samples of the validation split and bootstrap the paired difference over samples: last-round
minus per-round KV is \val{kvfull/x40/768/ced-enc/delta} [\val{kvfull/x40/768/ced-enc/delta/lo},
\val{kvfull/x40/768/ced-enc/delta/hi}] at 90M, \val{kvfull/x40/1024/ced-enc/delta}
[\val{kvfull/x40/1024/ced-enc/delta/lo}, \val{kvfull/x40/1024/ced-enc/delta/hi}] at 150M,
\val{kvfull/x40/1536/ced-enc/delta} [\val{kvfull/x40/1536/ced-enc/delta/lo}, \val{kvfull/x40/1536/ced-enc/delta/hi}]
at 340M, and \val{kvfull/long/1024/ced-enc/delta} [\val{kvfull/long/1024/ced-enc/delta/lo},
\val{kvfull/long/1024/ced-enc/delta/hi}] after \val{cfg/tokens/long}B tokens.

\paragraph{Downstream intervals.} CORE scores are mean centered accuracy over the suite's tasks. Their 95\%
intervals come from 1,000 bootstrap resamples of items, stratified by task; differences between two models are paired
over items where both have per-item results, and use unpaired resamples otherwise.

\section{Fixed-point analysis}
\label{app:fixedpoint}

\begin{figure}[t]
\centering
\safegraphic[width=\columnwidth]{figures/fig_fixedpoint.pdf}
\caption{\textbf{Fixed-point analysis.} The 150M CED encoder loop (\val{cfg/tokens/x40/1024}B tokens), trained with a
random or a fixed number of rounds and run for up to \val{cfg/fp/rounds} rounds on FineWeb-Edu validation data in fp32. (a) Median relative
change of the core state between rounds, $\rho_r = \lVert h_r - h_{r-1}\rVert / \lVert h_{r-1}\rVert$, over tokens,
with the interquartile band. (b) BPB after $r$ rounds minus BPB after the trained mean of \val{cfg/rounds} rounds
(dashed line).}
\label{fig:fixedpoint}
\end{figure}

\begin{figure*}[t]
\centering
\safegraphic[width=\textwidth]{figures/fig_longctx_corpora.pdf}
\caption{\textbf{Gap to dense on each corpus.} BPB minus the dense model with the same loop setting, by tokens of
context, for each corpus (150M
models, \val{cfg/tokens/x40/1024}B tokens). Markers are measurements; lines are least-squares fits in $\log_2$ of
the context beyond the SWA window (red dashed line). Top: no loop. Bottom: loop over layers 2--5. FineWeb-Edu counts
each token's context within its own document and is not averaged over offsets.}
\label{fig:longctx-corpora}
\end{figure*}

We run each loop for up to \val{cfg/fp/rounds} rounds and track $\rho_r$, the relative change of the core state
between rounds, as \citet{benice} do. With random rounds, $\rho_r$ shrinks by a factor of about
\val{fp/x40/1024/ced-enc/fw/contract} per round, from \val{fp/x40/1024/ced-enc/fw/rho/r4}\% at round 4 to
\val{fp/x40/1024/ced-enc/fw/rho/r8}\% at the \val{cfg/rounds} rounds we deploy and
\val{fp/x40/1024/ced-enc/fw/rho/r32}\% at round \val{cfg/fp/rounds}: the core contracts to a fixed point. Between
rounds 7 and 8, layer 2's global KV changes by only \val{fp/x40/1024/ced-enc/fw/kvstep/r8}\%, which is why every
round can read the last round's. From round 8 to \val{cfg/fp/rounds}, BPB stays within
\val{fp/x40/1024/ced-enc/fw/bpbspread/late} of its 8-round value, so running more rounds than trained is safe. The
same holds after \val{cfg/tokens/long}B tokens ($\rho_8 = \val{fp/long/1024/ced-enc/fw/rho/r8}\%$), for the decoder
loop ($\val{fp/x40/1024/ced-dec/fw/rho/r8}\%$) and for the 45M CED ($\val{fp/p140/768/p140-enc/fw/rho/r8}\%$). With
a fixed number of rounds, $\rho_r$ never falls below \val{fp/x40/1024/ced-enc-fixed/fw/rhomin/late}\% after round
8, and layer 2's KV changes by \val{fp/x40/1024/ced-enc-fixed/fw/kvstep/r8}\% per round. The state does not
converge but cycles with period \val{fp/x40/1024/ced-enc-fixed/fw/period}: BPB matches its 8-round value whenever
$r$ is a multiple of \val{fp/x40/1024/ced-enc-fixed/fw/period} and is
\val{fp/x40/1024/ced-enc-fixed/fw/bpbbad/min}--\val{fp/x40/1024/ced-enc-fixed/fw/bpbbad/max} higher otherwise. The
45M fixed-round CED shows the same cycle. Every fixed-round model we have backpropagates through only the last
\val{cfg/bwd} of its \val{cfg/rounds} rounds, which likely causes the period-\val{fp/x40/1024/ced-enc-fixed/fw/period}
cycle; we do not claim it for fixed-depth training in general. These models thus learn an orbit that reaches a
good state at the trained depth, not a fixed point, and the last round's KV is then a poor stand-in for earlier rounds'.

\section{Hyperparameters}
\label{app:hparams}

Every 12-layer run uses the hyperparameters of Table~\ref{tab:hparams}; the three sizes differ only in width, heads,
steps and training tokens.

\begin{table}[tb]
\caption{\textbf{Hyperparameters.} Shared by every 12-layer run, read from the run configurations.}
\label{tab:hparams}
\centering
\footnotesize
\setlength{\tabcolsep}{3pt}
\begin{tabular}{@{}l>{\raggedright\arraybackslash}p{0.5\columnwidth}@{}}
\toprule
Layers & \val{cfg/layers} \\
Model size (non-embedding) & 90M / 150M / 340M \\
Hidden size $d$ & 768 / 1024 / 1536 \\
Heads (query = KV) & \val{cfg/heads/768} / \val{cfg/heads/1024} / \val{cfg/heads/1536}, dimension \val{cfg/headdim} \\
FFN & \val{cfg/ffnmult}$d$, ReLU$^2$ \\
Value embeddings & layers \val{cfg/ve} \\
SWA window & \val{cfg/window} \\
RoPE $\theta$ & \val{cfg/rope} \\
Vocabulary & \val{cfg/vocab} (Qwen3) \\
Sequence length & \val{cfg/seq} \\
Batch & \val{cfg/batch} sequences \\
Optimizer & Muon (matrices), AdamW (others) \\
Learning rate & \val{cfg/lr}, warmup \val{cfg/warmup}\%, linear decay over the last \val{cfg/decay}\% \\
Weight decay & \val{cfg/wd} (AdamW), \val{cfg/muonwd} (Muon) \\
Steps & \val{cfg/steps/x40/768} / \val{cfg/steps/x40/1024} / \val{cfg/steps/x40/1536} \\
Tokens & \val{cfg/tokens/x40/768}B / \val{cfg/tokens/x40/1024}B / \val{cfg/tokens/x40/1536}B; \val{cfg/tokens/long}B for 150M \\
Rounds & $\max(1, \mathrm{Poisson}(\val{cfg/rounds}))$ per sequence; \val{cfg/rounds} at evaluation \\
Backpropagated rounds & last \val{cfg/bwd} \\
Loop-state initialization & \val{cfg/stateinit} (truncated normal); decay gate $\sqrt{\val{cfg/decayinit}}$ \\
\bottomrule
\end{tabular}
\end{table}

\section{Test-time rounds}
\label{app:rounds}

Table~\ref{tab:rounds} evaluates the 150M looped models with $T=1$ to 16 rounds. Running 12 or 16 rounds instead of
the trained mean of \val{cfg/rounds} lowers BPB by at most 0.0004 and never raises it, so these models can run more
rounds than they were trained with (Appendix~\ref{app:fixedpoint}).

\begin{table}[tb]
\caption{\textbf{Test-time rounds.} BPB when evaluating with $T$ rounds; the models were trained with \val{cfg/rounds} rounds on average.
Parallel forward pass with per-round KV on the first \val{cfg/evalwin/x40} validation samples, as in
Table~\ref{tab:kvshare}.}
\label{tab:rounds}
\centering
\footnotesize
\setlength{\tabcolsep}{2pt}
\resizebox{\columnwidth}{!}{\begin{tabular}{@{}lcccccc@{}}
\toprule
Model & $T{=}1$ & 2 & 4 & 8 & 12 & 16 \\
\midrule
CED, encoder loop & \sval{tt/x40/1024/ced-enc/T1} & \sval{tt/x40/1024/ced-enc/T2} & \sval{tt/x40/1024/ced-enc/T4} & \sval{tt/x40/1024/ced-enc/T8} & \sval{tt/x40/1024/ced-enc/T12} & \sval{tt/x40/1024/ced-enc/T16} \\
CED, enc.\ loop, \val{cfg/tokens/long}B & \sval{tt/long/1024/ced-enc/T1} & \sval{tt/long/1024/ced-enc/T2} & \sval{tt/long/1024/ced-enc/T4} & \sval{tt/long/1024/ced-enc/T8} & \sval{tt/long/1024/ced-enc/T12} & \sval{tt/long/1024/ced-enc/T16} \\
CED, decoder loop & \sval{tt/x40/1024/ced-dec/T1} & \sval{tt/x40/1024/ced-dec/T2} & \sval{tt/x40/1024/ced-dec/T4} & \sval{tt/x40/1024/ced-dec/T8} & \sval{tt/x40/1024/ced-dec/T12} & \sval{tt/x40/1024/ced-dec/T16} \\
Dense, loop L2--5 & \sval{tt/x40/1024/dense-l25/T1} & \sval{tt/x40/1024/dense-l25/T2} & \sval{tt/x40/1024/dense-l25/T4} & \sval{tt/x40/1024/dense-l25/T8} & \sval{tt/x40/1024/dense-l25/T12} & \sval{tt/x40/1024/dense-l25/T16} \\
Dense, loop L4--7 & \sval{tt/x40/1024/dense-l47/T1} & \sval{tt/x40/1024/dense-l47/T2} & \sval{tt/x40/1024/dense-l47/T4} & \sval{tt/x40/1024/dense-l47/T8} & \sval{tt/x40/1024/dense-l47/T12} & \sval{tt/x40/1024/dense-l47/T16} \\
\bottomrule
\end{tabular}}
\end{table}

\section{Fewer rounds for prompt tokens}
\label{app:split}

The encoder loop's prefill cost scales with the number of rounds that prompt tokens run. The unified evaluation
splits each validation sequence after its first half. It runs the first half with $T_p$ rounds and caches its
final-round KV, then scores the second half with $T_d$ rounds. Table~\ref{tab:split} shows BPB on the second
half only, so it is not comparable to the other tables. A prefill at $T_p$ rounds runs the same layers as a $T=T_p$
model. On $4\times\text{GB200}$, the measured prefill of $T=4$ and $T=2$ models is $2.72\times$ and $1.54\times$ that of
unlooped \ced, against $5.12\times$ at $T=8$ (Table~\ref{tab:unit-costs-app}).

\begin{table}[tb]
\caption{\textbf{Fewer rounds for prompt tokens.} CED encoder loop after \val{cfg/tokens/long}B tokens: BPB on the second half of each sequence, with
the first half (the prompt) run at $T_p$ rounds and the second half at $T_d$ rounds.}
\label{tab:split}
\centering
\small
\begin{tabular}{lcccc}
\toprule
$T_p$ & $T_d{=}2$ & 4 & 8 & 16 \\
\midrule
1 & \sval{split/long/1024/ced-enc/tp1_td2} & \sval{split/long/1024/ced-enc/tp1_td4} & \sval{split/long/1024/ced-enc/tp1_td8} & \sval{split/long/1024/ced-enc/tp1_td16} \\
2 & \sval{split/long/1024/ced-enc/tp2_td2} & \sval{split/long/1024/ced-enc/tp2_td4} & \sval{split/long/1024/ced-enc/tp2_td8} & \sval{split/long/1024/ced-enc/tp2_td16} \\
4 & \sval{split/long/1024/ced-enc/tp4_td2} & \sval{split/long/1024/ced-enc/tp4_td4} & \sval{split/long/1024/ced-enc/tp4_td8} & \sval{split/long/1024/ced-enc/tp4_td16} \\
8 & \sval{split/long/1024/ced-enc/tp8_td2} & \sval{split/long/1024/ced-enc/tp8_td4} & \sval{split/long/1024/ced-enc/tp8_td8} & \sval{split/long/1024/ced-enc/tp8_td16} \\
16 & \sval{split/long/1024/ced-enc/tp16_td2} & \sval{split/long/1024/ced-enc/tp16_td4} & \sval{split/long/1024/ced-enc/tp16_td8} & \sval{split/long/1024/ced-enc/tp16_td16} \\
\bottomrule
\end{tabular}
\end{table}

\begin{table}[tb]
\caption{\textbf{How fast BPB gaps grow with context.} The least-squares slope of BPB(a)$-$BPB(b) over the four context
buckets beyond the SWA window (128--255, \dots, 1024--2047 tokens) against $\log_2$ of the context, in
$10^{-3}$ BPB per doubling. Positive means that a falls further behind b as the context grows. Long docs pools PG19
test, FinePDFs and GovReport, with a 95\% paired bootstrap interval; FW is the FineWeb-Edu validation set. Hidden
size 1024, \val{cfg/tokens/x40/1024}B tokens.}
\label{tab:slope}
\centering
\small
\setlength{\tabcolsep}{2pt}
\begin{tabular}{lccccc}
\toprule
a $-$ b & Long docs & PG19 & PDFs & Gov. & FW \\
\midrule
\multicolumn{6}{l}{\emph{Gap to dense, loop L2--5}} \\
CED enc.\ loop & \sval{slopem/x40/1024/ced-enc-minus-dense-l25/longdocs} {\scriptsize[\sval{slopem/x40/1024/ced-enc-minus-dense-l25/longdocs/lo}, \sval{slopem/x40/1024/ced-enc-minus-dense-l25/longdocs/hi}]} & \sval{slopem/x40/1024/ced-enc-minus-dense-l25/pg19-test} & \sval{slopem/x40/1024/ced-enc-minus-dense-l25/finepdfs-long} & \sval{slopem/x40/1024/ced-enc-minus-dense-l25/govreport-test} & \sval{slopem/x40/1024/ced-enc-minus-dense-l25/fineweb-edu-100bt} \\
YOCO-U loop & \sval{slopem/x40/1024/yocou-loop-minus-dense-l25/longdocs} {\scriptsize[\sval{slopem/x40/1024/yocou-loop-minus-dense-l25/longdocs/lo}, \sval{slopem/x40/1024/yocou-loop-minus-dense-l25/longdocs/hi}]} & \sval{slopem/x40/1024/yocou-loop-minus-dense-l25/pg19-test} & \sval{slopem/x40/1024/yocou-loop-minus-dense-l25/finepdfs-long} & \sval{slopem/x40/1024/yocou-loop-minus-dense-l25/govreport-test} & \sval{slopem/x40/1024/yocou-loop-minus-dense-l25/fineweb-edu-100bt} \\
\multicolumn{6}{l}{\emph{Gap to dense, no loop}} \\
CED & \sval{slopem/x40/1024/ced-noloop-minus-dense-noloop/longdocs} {\scriptsize[\sval{slopem/x40/1024/ced-noloop-minus-dense-noloop/longdocs/lo}, \sval{slopem/x40/1024/ced-noloop-minus-dense-noloop/longdocs/hi}]} & \sval{slopem/x40/1024/ced-noloop-minus-dense-noloop/pg19-test} & \sval{slopem/x40/1024/ced-noloop-minus-dense-noloop/finepdfs-long} & \sval{slopem/x40/1024/ced-noloop-minus-dense-noloop/govreport-test} & \sval{slopem/x40/1024/ced-noloop-minus-dense-noloop/fineweb-edu-100bt} \\
YOCO-U & \sval{slopem/x40/1024/yocou-noloop-minus-dense-noloop/longdocs} {\scriptsize[\sval{slopem/x40/1024/yocou-noloop-minus-dense-noloop/longdocs/lo}, \sval{slopem/x40/1024/yocou-noloop-minus-dense-noloop/longdocs/hi}]} & \sval{slopem/x40/1024/yocou-noloop-minus-dense-noloop/pg19-test} & \sval{slopem/x40/1024/yocou-noloop-minus-dense-noloop/finepdfs-long} & \sval{slopem/x40/1024/yocou-noloop-minus-dense-noloop/govreport-test} & \sval{slopem/x40/1024/yocou-noloop-minus-dense-noloop/fineweb-edu-100bt} \\
\midrule
\multicolumn{6}{l}{\emph{Loop gain: no loop $-$ loop}} \\
Dense L2--5 & \sval{slopem/x40/1024/dense-noloop-minus-dense-l25/longdocs} {\scriptsize[\sval{slopem/x40/1024/dense-noloop-minus-dense-l25/longdocs/lo}, \sval{slopem/x40/1024/dense-noloop-minus-dense-l25/longdocs/hi}]} & \sval{slopem/x40/1024/dense-noloop-minus-dense-l25/pg19-test} & \sval{slopem/x40/1024/dense-noloop-minus-dense-l25/finepdfs-long} & \sval{slopem/x40/1024/dense-noloop-minus-dense-l25/govreport-test} & \sval{slopem/x40/1024/dense-noloop-minus-dense-l25/fineweb-edu-100bt} \\
CED encoder & \sval{slopem/x40/1024/ced-noloop-minus-ced-enc/longdocs} {\scriptsize[\sval{slopem/x40/1024/ced-noloop-minus-ced-enc/longdocs/lo}, \sval{slopem/x40/1024/ced-noloop-minus-ced-enc/longdocs/hi}]} & \sval{slopem/x40/1024/ced-noloop-minus-ced-enc/pg19-test} & \sval{slopem/x40/1024/ced-noloop-minus-ced-enc/finepdfs-long} & \sval{slopem/x40/1024/ced-noloop-minus-ced-enc/govreport-test} & \sval{slopem/x40/1024/ced-noloop-minus-ced-enc/fineweb-edu-100bt} \\
CED decoder & \sval{slopem/x40/1024/ced-noloop-minus-ced-dec/longdocs} {\scriptsize[\sval{slopem/x40/1024/ced-noloop-minus-ced-dec/longdocs/lo}, \sval{slopem/x40/1024/ced-noloop-minus-ced-dec/longdocs/hi}]} & \sval{slopem/x40/1024/ced-noloop-minus-ced-dec/pg19-test} & \sval{slopem/x40/1024/ced-noloop-minus-ced-dec/finepdfs-long} & \sval{slopem/x40/1024/ced-noloop-minus-ced-dec/govreport-test} & \sval{slopem/x40/1024/ced-noloop-minus-ced-dec/fineweb-edu-100bt} \\
YOCO-U & \sval{slopem/x40/1024/yocou-noloop-minus-yocou-loop/longdocs} {\scriptsize[\sval{slopem/x40/1024/yocou-noloop-minus-yocou-loop/longdocs/lo}, \sval{slopem/x40/1024/yocou-noloop-minus-yocou-loop/longdocs/hi}]} & \sval{slopem/x40/1024/yocou-noloop-minus-yocou-loop/pg19-test} & \sval{slopem/x40/1024/yocou-noloop-minus-yocou-loop/finepdfs-long} & \sval{slopem/x40/1024/yocou-noloop-minus-yocou-loop/govreport-test} & \sval{slopem/x40/1024/yocou-noloop-minus-yocou-loop/fineweb-edu-100bt} \\
\bottomrule
\end{tabular}
\end{table}

\begin{table*}[t]
\caption{\textbf{Long-document BPB by context bucket.} PG19 test, FinePDFs and GovReport (512, 512 and 396 samples
of 2048 tokens, each averaged over six offsets in its document); 150M models. One forward pass per sample with per-round KV and \val{cfg/rounds} rounds for looped
models. Rows above the line trained on \val{cfg/tokens/x40/1024}B tokens.}
\label{tab:longctx}
\centering
\small
\setlength{\tabcolsep}{4pt}
\begin{tabular}{lcccccc}
\toprule
Model & 1--127 & 128--255 & 256--511 & 512--1023 & 1024--2047 & mean \\
\midrule
Dense, no loop & \sval{pos/x40/1024/dense-noloop/longdocs/c1-128} & \sval{pos/x40/1024/dense-noloop/longdocs/c128-256} & \sval{pos/x40/1024/dense-noloop/longdocs/c256-512} & \sval{pos/x40/1024/dense-noloop/longdocs/c512-1024} & \sval{pos/x40/1024/dense-noloop/longdocs/c1024-2048} & \sval{pos/x40/1024/dense-noloop/longdocs/mean} \\
Dense, loop L2--5 & \sval{pos/x40/1024/dense-l25/longdocs/c1-128} & \sval{pos/x40/1024/dense-l25/longdocs/c128-256} & \sval{pos/x40/1024/dense-l25/longdocs/c256-512} & \sval{pos/x40/1024/dense-l25/longdocs/c512-1024} & \sval{pos/x40/1024/dense-l25/longdocs/c1024-2048} & \sval{pos/x40/1024/dense-l25/longdocs/mean} \\
Dense, loop L4--7 & \sval{pos/x40/1024/dense-l47/longdocs/c1-128} & \sval{pos/x40/1024/dense-l47/longdocs/c128-256} & \sval{pos/x40/1024/dense-l47/longdocs/c256-512} & \sval{pos/x40/1024/dense-l47/longdocs/c512-1024} & \sval{pos/x40/1024/dense-l47/longdocs/c1024-2048} & \sval{pos/x40/1024/dense-l47/longdocs/mean} \\
CED, no loop & \sval{pos/x40/1024/ced-noloop/longdocs/c1-128} & \sval{pos/x40/1024/ced-noloop/longdocs/c128-256} & \sval{pos/x40/1024/ced-noloop/longdocs/c256-512} & \sval{pos/x40/1024/ced-noloop/longdocs/c512-1024} & \sval{pos/x40/1024/ced-noloop/longdocs/c1024-2048} & \sval{pos/x40/1024/ced-noloop/longdocs/mean} \\
CED, encoder loop & \sval{pos/x40/1024/ced-enc/longdocs/c1-128} & \sval{pos/x40/1024/ced-enc/longdocs/c128-256} & \sval{pos/x40/1024/ced-enc/longdocs/c256-512} & \sval{pos/x40/1024/ced-enc/longdocs/c512-1024} & \sval{pos/x40/1024/ced-enc/longdocs/c1024-2048} & \sval{pos/x40/1024/ced-enc/longdocs/mean} \\
CED, decoder loop & \sval{pos/x40/1024/ced-dec/longdocs/c1-128} & \sval{pos/x40/1024/ced-dec/longdocs/c128-256} & \sval{pos/x40/1024/ced-dec/longdocs/c256-512} & \sval{pos/x40/1024/ced-dec/longdocs/c512-1024} & \sval{pos/x40/1024/ced-dec/longdocs/c1024-2048} & \sval{pos/x40/1024/ced-dec/longdocs/mean} \\
CED, enc.\ loop, fixed rounds & \sval{pos/x40/1024/ced-enc-fixed/longdocs/c1-128} & \sval{pos/x40/1024/ced-enc-fixed/longdocs/c128-256} & \sval{pos/x40/1024/ced-enc-fixed/longdocs/c256-512} & \sval{pos/x40/1024/ced-enc-fixed/longdocs/c512-1024} & \sval{pos/x40/1024/ced-enc-fixed/longdocs/c1024-2048} & \sval{pos/x40/1024/ced-enc-fixed/longdocs/mean} \\
YOCO-U style, no loop & \sval{pos/x40/1024/yocou-noloop/longdocs/c1-128} & \sval{pos/x40/1024/yocou-noloop/longdocs/c128-256} & \sval{pos/x40/1024/yocou-noloop/longdocs/c256-512} & \sval{pos/x40/1024/yocou-noloop/longdocs/c512-1024} & \sval{pos/x40/1024/yocou-noloop/longdocs/c1024-2048} & \sval{pos/x40/1024/yocou-noloop/longdocs/mean} \\
YOCO-U style, loop L2--5 & \sval{pos/x40/1024/yocou-loop/longdocs/c1-128} & \sval{pos/x40/1024/yocou-loop/longdocs/c128-256} & \sval{pos/x40/1024/yocou-loop/longdocs/c256-512} & \sval{pos/x40/1024/yocou-loop/longdocs/c512-1024} & \sval{pos/x40/1024/yocou-loop/longdocs/c1024-2048} & \sval{pos/x40/1024/yocou-loop/longdocs/mean} \\
\midrule
\multicolumn{7}{l}{\emph{\val{cfg/tokens/long}B training tokens}} \\
CED, no loop & \sval{pos/long/1024/ced-noloop/longdocs/c1-128} & \sval{pos/long/1024/ced-noloop/longdocs/c128-256} & \sval{pos/long/1024/ced-noloop/longdocs/c256-512} & \sval{pos/long/1024/ced-noloop/longdocs/c512-1024} & \sval{pos/long/1024/ced-noloop/longdocs/c1024-2048} & \sval{pos/long/1024/ced-noloop/longdocs/mean} \\
CED, encoder loop & \sval{pos/long/1024/ced-enc/longdocs/c1-128} & \sval{pos/long/1024/ced-enc/longdocs/c128-256} & \sval{pos/long/1024/ced-enc/longdocs/c256-512} & \sval{pos/long/1024/ced-enc/longdocs/c512-1024} & \sval{pos/long/1024/ced-enc/longdocs/c1024-2048} & \sval{pos/long/1024/ced-enc/longdocs/mean} \\
CED, decoder loop & \sval{pos/long/1024/ced-dec/longdocs/c1-128} & \sval{pos/long/1024/ced-dec/longdocs/c128-256} & \sval{pos/long/1024/ced-dec/longdocs/c256-512} & \sval{pos/long/1024/ced-dec/longdocs/c512-1024} & \sval{pos/long/1024/ced-dec/longdocs/c1024-2048} & \sval{pos/long/1024/ced-dec/longdocs/mean} \\
Dense, loop L4--7 & \sval{pos/long/1024/dense-l47/longdocs/c1-128} & \sval{pos/long/1024/dense-l47/longdocs/c128-256} & \sval{pos/long/1024/dense-l47/longdocs/c256-512} & \sval{pos/long/1024/dense-l47/longdocs/c512-1024} & \sval{pos/long/1024/dense-l47/longdocs/c1024-2048} & \sval{pos/long/1024/dense-l47/longdocs/mean} \\
\bottomrule
\end{tabular}
\end{table*}

\begin{figure*}[t]
\centering
\safegraphic[width=\textwidth]{figures/fig_recall.pdf}
\caption{\textbf{In-context recall.} BPB on the second half of samples whose tokens 1,024--2,047 repeat tokens
0--1,023 verbatim (PG19 test and arXiv, \val{cfg/recall/windows} samples), per 64-token bin; the first copied bin is
left out. Each panel compares the models of one size and token budget. Solid: looped (L2--5; decoder loop L7--10;
dense L4--7 in the last panel); dashed: unlooped.}
\label{fig:recall}
\end{figure*}

\section{Long-context evaluation}
\label{app:longctx}

\paragraph{Long documents.} PG19 test and FinePDFs contribute 512 samples and GovReport 396, each 2048 contiguous
tokens of one document starting at least 1,024 tokens in; every model scores the same samples. To keep the
difficulty of particular passages out of the position-wise curves, every sample is also scored after moving it
256, 512, 768 and 1,024 tokens earlier in its document and by one random offset of up to 1,024 tokens, and each
context bin averages the six copies, so the text at a given position comes from six places in the document. One forward pass scores each sample,
with \val{cfg/rounds} rounds and per-round KV for looped models. ``Long docs'' pools the three corpora.
Differences between models carry 95\% intervals from a paired bootstrap over samples (\val{cfg/bootreps}
resamples, stratified by corpus). For FineWeb-Edu we take each token's position in its own document as its
context.

\paragraph{BPB by context length.} Table~\ref{tab:longctx} gives the BPB behind Figure~\ref{fig:longctx} for every model, pooled over the three
long-document corpora, by the number of tokens of context. Figure~\ref{fig:longctx-corpora} repeats
Figure~\ref{fig:longctx}b--c for each corpus, and Table~\ref{tab:slope} gives how fast each gap grows with context.

\paragraph{In-context recall.}
\label{app:recall}
To test recall directly, we take the first 1,024 tokens of each PG19 test sample and of 512 arXiv documents from the Proof-Pile-2 test
split~\citep{llemma}, and repeat them verbatim at positions 1,024--2,047. Each copied token lies 1,024 tokens after its source, far beyond the SWA window, so only global KV can retrieve it.
Figure~\ref{fig:recall} shows BPB on the repeated half. Unlooped CED recalls
worse than unlooped dense (\val{recalldiff/x40/768/ced-noloop-minus-dense-noloop} and
\val{recalldiff/x40/1024/ced-noloop-minus-dense-noloop}); the encoder loop removes this deficit
(\val{recalldiff/x40/768/ced-enc-minus-ced-noloop} and \val{recalldiff/x40/1024/ced-enc-minus-ced-noloop} against
unlooped CED), whereas the decoder loop barely changes recall (\val{recalldiff/x40/768/ced-dec-minus-ced-noloop} and
\val{recalldiff/x40/1024/ced-dec-minus-ced-noloop}). The YOCO-U style loop, whose only global KV comes from layer 6,
recalls worst of the looped models (\val{recalldiff/x40/768/yocou-loop-minus-dense-l25} and
\val{recalldiff/x40/1024/yocou-loop-minus-dense-l25} against the dense loop). The differences shrink with scale: at
340M, all five models lie between \val{recall/x40/1536/dense-l25/bpb} and \val{recall/x40/1536/ced-noloop/bpb} BPB.
The CED encoder loop recalls better than the dense loop at 90M and 150M
(\val{recalldiff/x40/768/ced-enc-minus-dense-l25} and \val{recalldiff/x40/1024/ced-enc-minus-dense-l25} BPB), trails
it by \val{recalldiff/x40/1536/ced-enc-minus-dense-l25} [\val{recalldiff/x40/1536/ced-enc-minus-dense-l25/lo},
\val{recalldiff/x40/1536/ced-enc-minus-dense-l25/hi}] at 340M, and recalls better than the L4--7 dense loop of
Table~\ref{tab:long} after \val{cfg/tokens/long}B tokens (\val{recalldiff/long/1024/ced-enc-minus-dense-l47}). Across
the three sizes and both token budgets, it thus recalls within 0.03 BPB of the dense loop while caching one-sixth of its
global KV.


\section{Production Serving Setup and Detailed Results}
\label{app:serving}

\begin{table*}[t]
\centering
\caption{\textbf{Unit prefill and decode cost multipliers relative to unlooped \ced ($T=1$) across $4\times\text{GB200}$ (TP4) and $8\times\text{H200}$ (TP8).}
The analytic columns are a roofline model of FLOPs and bytes at GB200 peak rates, with expert parallelism over 4 GPUs and all of $L_{20}$ run on every prompt token, as in SGLang.
Prefill $a$ is the per-token slope of $t_{\text{pre}}(L) = b + aL + cL^2$, fitted over uncached prompts of 512--65{,}536 tokens at batch 1; for \ced $T=1$ it is $18.8\,\mu$s per prompt token per GB200 tray and $60.8\,\mu$s per H200 node. Decode is the per-step time averaged over 8K and 32K contexts at batches 16--256; for \ced $T=1$ the step at 8K context and batch 32 takes 8.6\,ms. GB200 120K uses distinct 120K-token contexts (per-layer KV, $T=8$: batches 1 and 8 only). Unit costs do not depend on the round policy, so each row holds for last-round and per-round KV alike.}
\label{tab:unit-costs-app}
\vskip 0.06in
\footnotesize
\setlength{\tabcolsep}{2.4pt}
\begin{tabular}{@{}lccccccc@{}}
\toprule
& \multicolumn{3}{c}{\textbf{Analytic Roofline Model}} & \multicolumn{3}{c}{\textbf{Measured $4\times\text{GB200}$ (TP4)}} & \textbf{Measured $8\times\text{H200}$ (TP8)} \\
\cmidrule(lr){2-4} \cmidrule(lr){5-7} \cmidrule(lr){8-8}
\textbf{Model Variant} & Prefill $a$ & Decode $8\text{--}32\text{K}$ & KV Pool & Prefill $a$ & Decode $8\text{--}32\text{K}$ & Decode $120\text{K}$ & Prefill $a$ \\
\midrule
\textbf{\ced, no loop (reference)} & $1.00\times$ & $1.00\times$ & $1.00\times$ & $1.00\times$ & $1.00\times$ & $1.00\times$ & $1.00\times$ \\
\textbf{Stock SGLang (no plugin)} & $1.00\times$ & $1.00\times$ & $1.00\times$ & $1.00\times$ & $1.02\times$ & --- & $0.98\times$ \\
\textbf{\ourmethod, $T=2$} & $1.87\times$ & $1.45\times$ & $1.00\times$ & $1.54\times$ & $1.39\times$ & --- & $1.61\times$ \\
\textbf{\ourmethod, $T=4$} & $3.60\times$ & $2.34\times$ & $1.00\times$ & $2.72\times$ & $2.13\times$ & --- & $2.87\times$ \\
\textbf{\ourmethod, $T=8$} & $7.06\times$ & $4.12\times$ & $\mathbf{1.00\times}$ & $\mathbf{5.12\times}$ & $\mathbf{3.95\times}$ & $\mathbf{3.83\times}$ & $\mathbf{5.42\times}$ \\
\textbf{\ced decoder loop, $T=8$} & $1.00\times^\ddagger$ & $4.13\times$ & $\mathbf{1.00\times}$ & $2.76\times^\ddagger$ & $\mathbf{3.99\times}$ & $\mathbf{3.89\times}$ & $1.80\times^\ddagger$ \\
\midrule
\textbf{Per-layer KV, no loop} & $1.91\times$ & $1.12\times$ & $0.086\times$ & --- & --- & --- & $1.73\times$ \\
\textbf{Per-layer KV, $T=8$} & $8.21\times$ & $4.23\times$ & $0.086\times$ & $\mathbf{6.97\times}$ & $\mathbf{4.24\times}$ & $\mathbf{4.56\times}$ & $\mathbf{6.88\times}$ \\
\bottomrule
\multicolumn{8}{@{}p{0.97\textwidth}@{}}{\scriptsize $^\ddagger$Analytically, the decoder loop runs its extra rounds only on the 128-token tail ($a\times 1.00$, $b\times 7.64$). Measured, its prefill slope is $2.76\times$ (GB200) and $1.80\times$ (H200), and on GB200 its per-request prefill time is $2.6\text{--}2.9\times$ that of \ced $T=1$ at every prompt length from 512 tokens (a single chunk) to 64K. The architecture does not require this extra cost; it comes from our implementation of the decoder-loop prefill.} \\
\end{tabular}
\end{table*}

\begin{table*}[t]
\centering
\caption{\textbf{Pure decode throughput saturation (last-round vs.\ per-round KV) for the \ced encoder loop} on $4\times\text{GB200}$ (real weights; each request has its own context, whose KV is allocated without being computed, and generates 256 tokens). Each cell is the ratio of the two policies' highest decode throughput; in parentheses, the batch at which each was measured. Last-round KV at 8K is capped at 768 requests (pool 80\% full), so the 8K ratios are lower bounds.}
\label{tab:hbm-decode-app}
\vskip 0.06in
\footnotesize
\setlength{\tabcolsep}{5pt}
\begin{tabular}{@{}lccccc@{}}
\toprule
\textbf{Rounds ($T$)} & \makecell{\textbf{Per-round}\\\textbf{KV factor}} & \makecell{\textbf{Pool (last-round}\\\textbf{/ per-round)}} & \textbf{ISL = 8K} & \textbf{ISL = 16K} & \textbf{ISL = 32K} \\
\midrule
\textbf{$T=2$} & $1.60\times$ & $8.11\text{M} \;/\; 5.07\text{M}$ & $\mathbf{\ge 1.12\times}$ ($768$ vs.\ $570$) & $\mathbf{1.17\times}$ ($487$ vs.\ $288$) & $\mathbf{1.34\times}$ ($245$ vs.\ $144$) \\
\textbf{$T=4$} & $2.80\times$ & $8.11\text{M} \;/\; 2.90\text{M}$ & $\mathbf{\ge 1.42\times}$ ($768$ vs.\ $342$) & $\mathbf{1.58\times}$ ($487$ vs.\ $174$) & $\mathbf{1.80\times}$ ($245$ vs.\ $87$) \\
\textbf{$T=8$} & $5.20\times$ & $8.11\text{M} \;/\; 1.56\text{M}$ & $\mathbf{\ge 2.05\times}$ ($768$ vs.\ $184$) & $\mathbf{2.35\times}$ ($487$ vs.\ $93$) & $\mathbf{2.83\times}$ ($245$ vs.\ $47$) \\
\bottomrule
\end{tabular}
\end{table*}

\paragraph{DeepSeek-V4.1-Flash architecture and SGLang plugin.}
DeepSeek-V4.1-Flash~\citep{deepseekv41} is a 40-layer Causal Encoder--Decoder MoE model (hidden dimension $5{,}120$, 384 routed experts per MoE layer with top-6 routing, FP8 attention/dense projections, and MXFP4 expert weights) partitioned into a 21-layer encoder (\Lr{0}{20}) and a 19-layer decoder (\Lr{21}{39}).
Every layer maintains a 128-token sliding-window attention (SWA) ring buffer ($768$ reserved slots per active request in SGLang, totaling $17.94$\,MB/request).
Long-range context is served by only four global KV-source layers ($L_2, L_8, L_{14}, L_{20}$) with compression ratios $(2, 2, 2, 1)$: the three sources inside the looped range \Lr{2}{19} store $3 \times 1/2 = 1.5$ compressed entries per token, and $L_{20}$, whose KV the whole decoder reads, stores one ($2.5$ compressed entries/token total).
In SGLang (\texttt{da64c5cbb}), each compressed entry stores $584$\,B of MLA latent KV (\texttt{fp8\_e4m3} nope + BF16 RoPE + scale) plus $68$\,B of FP4 lightning-indexer keys ($652$\,B/entry), giving $2.5 \times 652 = 1{,}630$\,B/token for \ced vs.\ $29.0 \times 652 = 18{,}908$\,B/token for the per-layer-KV baseline, in which all 38 layers \Lr{2}{39} own distinct global KV caches ($11.60\times$ larger; the analytical FP4-latent cost model uses $288 + 68 = 356$\,B/entry, or $890$\,B vs.\ $10{,}324$\,B/token, with identical ratios).
Under per-round KV, looping \Lr{2}{19} for $T$ rounds multiplies global KV by $1 + \frac{1.5}{2.5}(T-1) = 1 + 0.6(T-1)$ for \ced ($5.20\times$ at $T=8$, or $8{,}476$\,B/tok) and by $1 + \frac{9}{29}(T-1)$ with per-layer KV ($3.17\times$ at $T=8$, or $59{,}984$\,B/tok).

Our weight-tied plugin (\texttt{sglang\_ced\_loop}) wraps SGLang's layer execution loop into a virtual schedule of $(l, r)$ passes while keeping a single copy of physical weights.
For loops whose first layer reuses an earlier source's sparse top-$k$ indices (the decoder loop \Lr{21}{38} reusing $L_{20}$), the plugin snapshots the top-$k$ buffer on round 1 and restores it on rounds $2\dots T$; for $L_{14}$'s Engram table, the embedding lookup and key-value projection execute once on round 1 while the hidden-state gate runs every round.
At $T=1$, the plugin's unit prefill and decode costs are within 2\% of stock SGLang on both $8\times\text{H200}$ (TP8/EP8) and $4\times\text{GB200}$ (TP4/EP4) (Table~\ref{tab:unit-costs-app}), with zero layer-count or NLL discrepancies.
Table~\ref{tab:unit-costs-app} reports the complete unit prefill and decode cost multipliers, and Table~\ref{tab:hbm-decode-app} details the HBM-saturated pure decode sweep across context lengths and rounds.


\paragraph{Setup of Figure~\ref{fig:p1-pareto}.}
\label{app:serving-fig}
All runs use one $4\times\text{GB200}$ tray (TP4/EP4) and one SGLang server.
Within each panel, all models share the server settings, apart from the pool cap that emulates per-round KV. The settings differ between experiments, and so does the KV pool of \ourmethod: $5.11$\,M tokens in panels (a) and (b) and in Table~\ref{tab:cost}, $7.9$\,M in panel (c), and $8.11$\,M in Table~\ref{tab:hbm-decode-app}.
Panels (a) and (b) use the \vfour-Flash architecture with dummy weights, which change the ratios between models by 2--4\% (Appendix~\ref{app:serving-more}); panel (c) uses the real checkpoint weights.
Per-round KV is emulated with identical compute: the KV pool is capped by the per-round byte factor ($5.2\times$ for \ced at $T=8$, $3.17\times$ with per-layer KV; Appendix~\ref{app:serving}).
The baseline with per-layer KV gives every layer \Lr{2}{39} its own global KV and runs every layer on every prompt token in prefill (full replay).
A point is \emph{preemption-free} when no request is retracted; panel (a) draws only such points, and the dotted line carries each baseline's last preemption-free throughput to $c=512$.
\emph{(a)}~8K/1K requests are random token ids with \texttt{ignore\_eos} and $\max(20, 5c)$ prompts per point (one or two repeats).
The baseline (per-layer and per-round KV) fills its pool at $c=16$ ($T=4$) and $c=8$ ($T=8$), \ced with per-round KV at $c=128$ and $c=64$; last-round KV stays preemption-free up to $c=512$.
At $c=512$, last-round KV gives $1.27\times$ ($T=4$) and $1.30\times$ ($T=8$) the peak throughput of per-round KV, and $3.1\times$ and $3.6\times$ that of the baseline.
\emph{(b)}~The 128K groups are pure decode: each request has its own 131K-token context and generates 256 tokens, at the largest batch the KV pool holds ($38$ requests with last-round KV; $13$ and $7$ with per-round KV at $T=4$ and $8$; $1$ for the baseline at $T=4$ and none at $T=8$).
There, last-round KV gives $1.9\times$ ($T=4$) and $2.9\times$ ($T=8$) the throughput of per-round KV.
The ShareGPT group compares \ourmethod with the baseline only.
\emph{(c)}~Mooncake \texttt{toolagent} trace~\citep{mooncake}, closed loop with $c$ concurrent sessions, encoder loop at $T=8$.
With per-round KV the prefix cache holds $1.5$\,M tokens and is full from $c=128$, so evicted session histories are recomputed through all eight rounds; last-round KV holds $7.9$\,M tokens and peaks at $77\%$ use at $c=512$.
Median TTFT drops $13.6\times$ at $c=256$ and $29.7\times$ at $c=512$, while output throughput is within $13\%$ between the two policies.

